% === START SEGMENT 2 ===

%% ----------------------------------------------------------------
%%  SECTION 6 — COMPARISON WITH EXISTING PLATFORMS
%% ----------------------------------------------------------------
\section{Comparison with Suricata, Snort, and Commercial Platforms}
\label{sec:comparison}

Table~\ref{tab:platform_comparison} extends the original three-system
comparison to include the major AI-augmented commercial security platforms
released since 2023.  We evaluate across twelve axes spanning detection
methodology, AI integration depth, and operational compliance.

\begin{table*}[t]
\centering
\caption{Feature comparison of RobustIDPS.ai against open-source IDS engines and commercial AI-augmented platforms.}
\label{tab:platform_comparison}
\small
\begin{tabular}{@{}lcccccccc@{}}
\toprule
\textbf{Capability} & \rotatebox{70}{\textbf{RobustIDPS.ai}} & \rotatebox{70}{\textbf{Suricata}} & \rotatebox{70}{\textbf{Snort 3}} & \rotatebox{70}{\textbf{CrowdStrike Charlotte AI}} & \rotatebox{70}{\textbf{SentinelOne Purple AI}} & \rotatebox{70}{\textbf{MS Security Copilot}} & \rotatebox{70}{\textbf{Google Sec-Gemini}} \\
\midrule
Signature-based detection       & \checkmark & \checkmark & \checkmark & \checkmark & \checkmark & \checkmark & \checkmark \\
ML-based detection (13+ models) & \checkmark & ---        & ---        & Partial    & Partial    & ---        & Partial    \\
Adversarial robustness testing  & \checkmark & ---        & ---        & ---        & ---        & ---        & ---        \\
Continual learning / EWC        & \checkmark & ---        & ---        & ---        & ---        & ---        & ---        \\
Constrained RL response         & \checkmark & ---        & ---        & ---        & ---        & ---        & ---        \\
LLM integration (multi-provider)& \checkmark & ---        & ---        & Proprietary& Proprietary& Proprietary& Proprietary\\
Agentic auto-investigation      & \checkmark & ---        & ---        & \checkmark & \checkmark & \checkmark & \checkmark \\
NL threat hunting               & \checkmark & ---        & ---        & \checkmark & \checkmark & \checkmark & \checkmark \\
LLM attack surface testing      & \checkmark & ---        & ---        & ---        & ---        & ---        & ---        \\
MITRE ATT\&CK mapping           & \checkmark & Partial    & Partial    & \checkmark & \checkmark & \checkmark & \checkmark \\
OWASP LLM Top 10 compliance    & \checkmark & ---        & ---        & ---        & ---        & Partial    & ---        \\
NIST AI RMF dashboard           & \checkmark & ---        & ---        & ---        & ---        & ---        & ---        \\
Post-quantum IDS module         & \checkmark & ---        & ---        & ---        & ---        & ---        & ---        \\
Open-source / auditable         & \checkmark & \checkmark & \checkmark & ---        & ---        & ---        & ---        \\
\bottomrule
\end{tabular}
\end{table*}

A key differentiator is that RobustIDPS.ai provides \emph{open, multi-provider}
LLM integration (Claude, GPT-4o, Gemini, DeepSeek, Llama) rather than a single
proprietary model, enabling SOC teams to benchmark reasoning quality across
providers.  Furthermore, no commercial platform currently offers integrated LLM
attack surface testing---a capability essential for organisations deploying
copilot-augmented SOC workflows.

%% ----------------------------------------------------------------
%%  SECTION 7 — SOC INTELLIGENCE SUITE
%% ----------------------------------------------------------------
\section{SOC Intelligence Suite}
\label{sec:soc_intelligence}

The SOC Intelligence Suite comprises six purpose-built pages that transform
raw detection output into actionable intelligence for Tier~1--3 analysts.

\subsection{Auto-Investigation}
\label{sec:auto_investigation}

The agentic pipeline executes four phases: (1)~\textbf{Evidence Collection}---retrieves the triggering flow, $\pm$30\,s temporal context, and correlated alerts; (2)~\textbf{Contextual Enrichment}---queries IP reputation, geo-location, and alert frequency; (3)~\textbf{LLM Reasoning}---submits enriched evidence to the selected LLM for root-cause analysis and MITRE ATT\&CK mapping; (4)~\textbf{Verdict Synthesis}---aggregates LLM output with model confidence into a formatted investigation report.

\subsection{Natural Language Threat Hunt}
\label{sec:nl_threat_hunt}

Analysts issue free-text queries (e.g., \texttt{"high-confidence DDoS on port 443 from Russian IPs"}).  The parser extracts six filter dimensions---\texttt{attack\_type}, \texttt{confidence}, \texttt{port}, \texttt{severity}, \texttt{protocol}, \texttt{destination}---and applies them against the detection log.

\subsection{Incident Report Generation}
\label{sec:incident_report}

Generates auto-formatted reports with executive summary, technical details, MITRE ATT\&CK coverage matrix, event timeline, and remediation steps.  Exportable as PDF and structured JSON for SIEM ingestion.

\subsection{Threat Intelligence}
\label{sec:threat_intel}

Per-IP enrichment: reputation score (0--100), ASN, geo-location with map visualisation, historical threat activity, and composite threat score.

\subsection{IDS Rule Generator}
\label{sec:rule_generator}

Translates detected attacks into deployable rules using 21 Suricata/Snort templates covering reconnaissance, DoS, exploitation, lateral movement, and exfiltration.  Output includes SID, priority, and reference CVEs.

\subsection{CVE Vulnerability Mapper}
\label{sec:cve_mapper}

Maps 30 attack classes to CVE identifiers via 10 curated mapping tables, each with CVSS score, affected versions, patches, and mitigations.

%% ----------------------------------------------------------------
%%  SECTION 8 — LLM SECURITY TESTING
%% ----------------------------------------------------------------
\section{LLM Security Testing}
\label{sec:llm_security}

The LLM Attack Surface module provides five interactive testing environments
for evaluating the robustness of LLM-integrated security workflows.
Table~\ref{tab:llm_detection} summarises detection rates across all modules.

\begin{table}[t]
\centering
\caption{LLM attack detection rates by module (DefensePipeline v2).}
\label{tab:llm_detection}
\begin{tabular}{@{}lcc@{}}
\toprule
\textbf{Attack Module} & \textbf{Categories} & \textbf{Detection \%} \\
\midrule
Prompt Injection        & 8  & 96.3 \\
Jailbreak Taxonomy      & 8  & 95.1 \\
RAG Poisoning           & 4  & 93.7 \\
Multi-Agent Chain       & 5  & 91.8 \\
Data Poisoning          & 4  & 95.6 \\
\midrule
\textbf{Overall (weighted)} & \textbf{29} & \textbf{94.5} \\
\bottomrule
\end{tabular}
\end{table}

\subsection{Prompt Injection Playground}
Eight injection categories are tested against a real \texttt{DefensePipeline}
instance: direct override, indirect injection via log fields, context window
manipulation, few-shot poisoning, encoding-based bypass, role-play escalation,
multi-turn progressive compromise, and tool-invocation hijacking.  Each test
returns the sanitised output, detection flag, matched rule, and latency.

\subsection{Jailbreak Taxonomy}
Eight jailbreak techniques are catalogued and testable in real time:
DAN-style persona override, base64 encoding bypass, token smuggling,
multi-language confusion, hypothetical framing, system prompt extraction,
payload splitting, and recursive self-improvement prompts.

\subsection{RAG Poisoning Simulator}
Simulates four poisoning vectors against the platform's RAG pipeline:
document injection, embedding space perturbation, metadata manipulation,
and chunk boundary exploitation.  Defence mechanisms include document
provenance verification, embedding similarity thresholds ($\tau = 0.82$),
and semantic consistency cross-checks.

\subsection{Multi-Agent Chain Attack Simulator}
Models a four-agent topology (Traffic Analyst, Enrichment Agent, Decision
Agent, Response Agent) with configurable inter-agent trust scores.  Five
attack patterns are simulated: trust delegation abuse, capability escalation
chains, consensus manipulation, memory poisoning propagation, and
orchestrator compromise.

\subsection{Data Poisoning Simulator}
Tests four poisoning strategies against the training pipeline: label flipping
(random and targeted), backdoor trigger injection, gradient-based poisoning,
and clean-label attacks.  Defence evaluation includes spectral signature
detection, activation clustering, and neural cleanse verification.

%% ----------------------------------------------------------------
%%  SECTION 9 — MLSECOPS STANDARDS COMPLIANCE
%% ----------------------------------------------------------------
\section{MLSecOps Standards Compliance}
\label{sec:mlsecops}

RobustIDPS.ai integrates three major security and AI governance frameworks
directly into the analyst workflow.

\subsection{MITRE ATT\&CK Mapper}
Maps all 30 detection classes to corresponding MITRE ATT\&CK techniques
across the full kill chain.  An interactive kill chain visualisation
highlights coverage gaps.  Each mapping includes technique ID, tactic,
detection logic reference, and confidence level.

\subsection{OWASP LLM Top 10 Scorecard}
Evaluates the platform against each of the OWASP Top 10 for LLM Applications
(2025 edition), providing a compliance scorecard with current mitigation
status, residual risk rating, and remediation recommendations for each item.

\subsection{NIST AI Risk Management Framework Dashboard}
Presents compliance status across the four NIST AI RMF functions (Govern,
Map, Measure, Manage) with per-subcategory evidence mapping drawn from
platform features, testing results, and operational telemetry.

\subsection{Alert Triage Classifier}
An ML-assisted triage system that classifies incoming alerts into severity
tiers (Critical, High, Medium, Low, Informational) based on attack class,
confidence, affected asset criticality, and historical false-positive rates.

%% ----------------------------------------------------------------
%%  SECTION 10 — MULTI-AGENT PQC-IDS
%% ----------------------------------------------------------------
\section{Multi-Agent PQC-IDS}
\label{sec:pqc_ids}

The Multi-Agent Post-Quantum Cryptography IDS module represents the
successor to the PQ-IDPS system developed in Branch~3, redesigned around
a cooperative multi-agent architecture for post-quantum traffic analysis.

\subsection{Architecture}
Four specialised agents collaborate through attention-weighted fusion:
\begin{itemize}
  \item \textbf{Traffic Analyst} --- Extracts flow-level features from
        PQC handshake traffic (key exchange sizes, cipher negotiation
        patterns, certificate chain lengths).
  \item \textbf{PQC Specialist} --- Classifies cryptographic algorithm
        families (CRYSTALS-Kyber, CRYSTALS-Dilithium, SPHINCS+, BIKE)
        and detects anomalous parameter choices.
  \item \textbf{Anomaly Detector} --- Applies statistical deviation
        analysis on flow timing, payload entropy, and session metadata.
  \item \textbf{Coordinator} --- Fuses agent outputs via learned
        attention weights and produces final classification with
        per-agent contribution scores.
\end{itemize}

The complete model comprises 70,598 trainable parameters with
attention-weighted fusion across all four agent outputs.

\subsection{PQC Traffic Lab}
The PQC Traffic Lab provides a three-dataset comparison environment
for evaluating detection performance across heterogeneous PQC traffic
captures.  Validation datasets are sourced from Kaggle
(DOI:~\texttt{10.34740/kaggle/dsv/15424420}) and include both synthetic
and real-world PQC session recordings.

%% ----------------------------------------------------------------
%%  SECTION 11 — SECURITY, PERFORMANCE, AND USER MANAGEMENT
%% ----------------------------------------------------------------
\section{Security, Performance, and User Management}
\label{sec:security_performance}

\subsection{Security Hardening}
All user-supplied content is sanitised through DOMPurify before rendering.
WebSocket connections are authenticated via JWT tokens with per-job
authorisation scoping.  HTTP responses include comprehensive security
headers: \texttt{Content-Security-Policy}, \texttt{X-Frame-Options},
\texttt{X-Content-Type-Options}, \texttt{Strict-Transport-Security},
and \texttt{Referrer-Policy}.

\subsection{Performance Optimisations}
A \texttt{fast\_mode} toggle reduces inference latency by disabling Monte
Carlo dropout passes and uncertainty quantification.  Batch processing uses
a configurable batch size (default 4096 flows) with context-aware
\texttt{MAX\_ROWS} limits that adapt to available system memory.  Detection
throughput exceeds 12,000 flows/second on a single CPU core in fast mode.

\subsection{User Profiles}
Each user maintains a profile comprising a unique RobustID, preferred LLM
provider and model, timezone, display preferences, and optional ORCID
identifier for research attribution.

\subsection{Admin Dashboard}
The administration interface provides four management tabs:
\begin{itemize}
  \item \textbf{Users} --- Account management, role assignment, API key
        provisioning.
  \item \textbf{Audit} --- Comprehensive activity log with filterable
        event types, timestamps, and actor identification.
  \item \textbf{Sessions} --- Active session monitoring with idle timeout
        enforcement and forced termination capability.
  \item \textbf{Health} --- System health metrics including model load
        status, memory utilisation, WebSocket connection count, and
        queue depths.
\end{itemize}

\subsection{NoticeBoard}
A platform-wide notification system with auto-dismiss timers, severity-based
colour coding, and persistent storage for compliance-relevant announcements.

%% ----------------------------------------------------------------
%%  SECTION 12 — VALIDATION DATASETS
%% ----------------------------------------------------------------
\section{Validation Datasets}
\label{sec:datasets}

\subsection{Crafted PCAP Collections}
Three purpose-built PCAP datasets provide controlled ground truth for
regression testing:
\begin{itemize}
  \item \textbf{Dataset A} --- 60,000 flows covering all 34 attack classes
        with balanced class representation.
  \item \textbf{Dataset B} --- 40,000 flows emphasising low-prevalence
        attack types (Heartbleed, Infiltration, SSH-Patator) for
        minority-class evaluation.
  \item \textbf{Dataset C} --- 50,000 flows with injected concept drift
        at three temporal boundaries for continual learning validation.
\end{itemize}

\subsection{PQC Datasets}
Post-quantum cryptography traffic datasets are sourced from Kaggle
(DOI:~\texttt{10.34740/kaggle/dsv/15424420}) and include CRYSTALS-Kyber,
CRYSTALS-Dilithium, and hybrid TLS~1.3 handshake captures.

\subsection{Ground Truth Format}
All datasets are accompanied by ground truth CSV files with columns:
\texttt{flow\_id}, \texttt{timestamp}, \texttt{src\_ip}, \texttt{dst\_ip},
\texttt{src\_port}, \texttt{dst\_port}, \texttt{protocol}, \texttt{label},
\texttt{attack\_class}, and \texttt{confidence}.

\subsection{Validation Methodology}
Each model is evaluated using stratified 5-fold cross-validation with
class-weighted metrics.  Adversarial robustness is assessed at six
perturbation budgets ($\epsilon \in \{0.005, 0.01, 0.02, 0.03, 0.05, 0.1\}$).
Continual learning retention is measured across five sequential drift episodes
using backward transfer (BWT) and forward transfer (FWT) metrics.

%% ----------------------------------------------------------------
%%  SECTION 13 — FUTURE DIRECTIONS
%% ----------------------------------------------------------------
\section{Future Directions}
\label{sec:future}

With Phase~1 (core platform, 46 pages, 13+ models, SOC Intelligence Suite,
LLM Security Testing, MLSecOps compliance) now complete, the roadmap
focuses on the following Phase~2 objectives:

\begin{enumerate}
  \item \textbf{GPU-accelerated inference} --- CUDA and TensorRT
        integration for sub-millisecond per-flow classification at
        100\,Gbps line rates.
  \item \textbf{Real multi-agent LLM orchestration} --- Deploying
        production multi-agent workflows where specialised LLM agents
        (triage, enrichment, response) coordinate through verified
        message-passing protocols.
  \item \textbf{Online continual learning with streaming} --- Replacing
        batch-mode EWC with streaming-compatible algorithms (e.g.,
        Online-EWC, PackNet) for real-time model adaptation.
  \item \textbf{Multi-modal fusion} --- Integrating network flow data
        with system call logs, endpoint telemetry, and DNS query streams
        for holistic threat detection.
  \item \textbf{Post-quantum IDS on constrained IoT} --- Deploying
        quantised Multi-Agent PQC-IDS models on ARM Cortex-M and
        RISC-V microcontrollers for edge-native PQC monitoring.
  \item \textbf{Formal verification of LLM defence pipelines} ---
        Applying SMT-based verification to prove bounded safety
        properties of the DefensePipeline under adversarial input.
  \item \textbf{Cross-platform benchmarking vs commercial SOC} ---
        Establishing reproducible benchmarks comparing RobustIDPS.ai
        against CrowdStrike Charlotte AI, SentinelOne Purple AI,
        Microsoft Security Copilot, and Google Sec-Gemini on
        standardised alert datasets.
\end{enumerate}

%% ----------------------------------------------------------------
%%  SECTION 14 — CONCLUSION
%% ----------------------------------------------------------------
\section{Conclusion}
\label{sec:conclusion}

RobustIDPS.ai has matured from a research prototype into a comprehensive,
production-grade AI-native intrusion detection and prevention platform
spanning 46 interconnected pages and integrating more than 13 machine
learning and deep learning models.  The platform uniquely combines
adversarially robust network intrusion detection (seven-branch surrogate
ensemble with EWC continual learning and CPO autonomous response),
a full SOC Intelligence Suite (agentic auto-investigation, natural language
threat hunting, automated incident reporting, threat intelligence
enrichment, IDS rule generation, and CVE vulnerability mapping),
systematic LLM attack surface testing (prompt injection, jailbreak,
RAG poisoning, multi-agent chain attacks, and data poisoning with a
94.5\% overall detection rate), MLSecOps compliance dashboards (MITRE
ATT\&CK, OWASP LLM Top~10, NIST AI RMF), and a Multi-Agent PQC-IDS
module for post-quantum traffic analysis.

By providing open-source, multi-provider LLM integration alongside
rigorous adversarial testing capabilities, RobustIDPS.ai addresses
a critical gap in the security ecosystem: enabling organisations to
deploy AI-augmented SOC workflows while systematically evaluating and
hardening the AI components themselves.  The platform serves both as
an operational security tool and as a research instrument for advancing
the state of the art in robust, trustworthy AI for cybersecurity.

%% ----------------------------------------------------------------
%%  SOFTWARE AVAILABILITY
%% ----------------------------------------------------------------
\section*{Software Availability}

\begin{itemize}
  \item \textbf{Platform:} \url{https://robustidps.ai}
  \item \textbf{Source Code:} \url{https://github.com/rogerpanel/robustidps.ai}
  \item \textbf{Zenodo Archive:} DOI:~\texttt{10.5281/zenodo.19129512}
  \item \textbf{Multi-Agent PQC-IDS:} \url{https://github.com/rogerpanel/Multi-Agent-PQC-models}
  \item \textbf{PQC Datasets:} DOI:~\texttt{10.34740/kaggle/dsv/15424420}
  \item \textbf{Contact:} \href{mailto:support@robustidps.ai}{support@robustidps.ai}
\end{itemize}

\end{document}
